\section{Neurology 101 --- What the Reader Actually Needs to Remember}\label{06_neurology_101-neurology-101-what-the-reader-actually-needs-to-remember}

\subsection{\texorpdfstring{A Chapter Dossier for \emph{I Have No Dopamine, and I Must Function}}{A Chapter Dossier for I Have No Dopamine, and I Must Function}}\label{06_neurology_101-a-chapter-dossier-for-i-have-no-dopamine-and-i-must-function}

\begin{center}\rule{0.5\linewidth}{0.5pt}\end{center}

\subsubsection{Opening note to myself}\label{06_neurology_101-opening-note-to-myself}

Last chapter, we crawled out of the Precambrian on the back of a lamprey. Fine. But the rest of this book is going to ask the reader to think about very specific failures --- D2/D3 downregulation, anhedonia, cognitive bradykinesia, the moment when ``just try harder'' becomes a sick joke --- and none of that lands without a working picture of which parts of the brain are doing what to whom.

So this chapter is the payoff. It is small. There are exactly six structures (or structure-clusters) you have to keep in your head from here on out. Everything later is a perturbation of the circuit they form together. If you can sketch that circuit on a napkin by the end, we are in business.

I am going to keep saying the same thing in different ways, because that is how this stuff sticks.

\begin{center}\rule{0.5\linewidth}{0.5pt}\end{center}

\subsection{The single picture you need}\label{06_neurology_101-the-single-picture-you-need}

Before any anatomy: there is one loop. The neuroscientific term is the \textbf{cortico-striato-thalamo-cortical loop}, sometimes shortened to CSTC. Five letters, one idea.

Cortex (especially prefrontal cortex) projects down into the \textbf{striatum}. The striatum talks to the rest of the \textbf{basal ganglia}, which sits there with its foot on the brake of every motor and cognitive program you have. The basal ganglia output gates the \textbf{thalamus}. The thalamus relays back up to cortex. Cortex now does the next thing --- which is itself another input to the striatum. Loop.

Two modulatory systems whisper into this loop and change how it behaves: the \textbf{dopamine} projections from the midbrain (VTA and substantia nigra), and the \textbf{limbic} inputs (amygdala, hippocampus, hypothalamus) that tell the loop what matters and why.

That is the whole architecture. The rest of this chapter is just unpacking the boxes.

\begin{center}\rule{0.5\linewidth}{0.5pt}\end{center}

\subsection{1. The Basal Ganglia --- the ancient action-selection engine}\label{06_neurology_101-the-basal-ganglia-the-ancient-action-selection-engine}

\textbf{Where it sits.} Buried under the cortex, deep in each hemisphere. A cluster of subcortical nuclei: striatum (caudate, putamen, nucleus accumbens), globus pallidus (internal and external segments, GPi and GPe), subthalamic nucleus (STN), and substantia nigra (with a pars compacta that makes dopamine and a pars reticulata that is, functionally, an output station like GPi).

\textbf{What it does.} Action selection. At any moment dozens of possible motor and cognitive programs are sitting in your cortex with their hands raised --- \emph{speak, stand up, scratch, ruminate about that email, reach for the cup}. The basal ganglia's job is to pick one and unmask it while keeping the others suppressed. By default it tonically inhibits the thalamus and brainstem motor centers, which is to say its resting state is ``no, don't do anything yet.'' \href{https://pmc.ncbi.nlm.nih.gov/articles/PMC3853485/}{PubMed Central} Selecting an action means \emph{releasing} that inhibition for one specific channel.

This is implemented through two parallel pathways from striatum:

\begin{itemize}
\tightlist
\item
  \textbf{Direct pathway (``go'').} Striatal medium spiny neurons (MSNs) expressing \textbf{D1 receptors} project to GPi/SNr and inhibit the inhibition. Net effect: thalamus is dis-inhibited, the chosen program runs. \href{https://pmc.ncbi.nlm.nih.gov/articles/PMC3853485/}{PubMed Central}
\item
  \textbf{Indirect pathway (``no-go'').} Striatal MSNs expressing \textbf{D2 receptors} route through GPe and STN to \emph{increase} GPi/SNr output. \href{https://pmc.ncbi.nlm.nih.gov/articles/PMC3853485/}{PubMed Central} Net effect: more inhibition, the program is held back.
\end{itemize}

Dopamine tilts this competition. D1 neurons are \emph{excited} by dopamine; D2 neurons are \emph{inhibited} by it. So a dopamine pulse simultaneously says ``yes'' to go and ``shh'' to no-go for whichever program is currently being addressed. Conversely, when dopamine is low or the dopamine signal is blunted, the indirect (no-go) pathway dominates and everything feels effortful, slow, and under-selected. That last sentence is the seed of everything I will later say about cognitive bradykinesia and the failure of ``willpower.''

\textbf{Connection to the book's argument.} This is the structure that decides whether an action happens. Agency, in the only sense that survives contact with neuroscience, lives here. When people speak of ``willpower'' they almost always mean some property of cortex bullying basal ganglia into selecting the harder action over the easier one. But cortex does not select; it \emph{votes}. The actual selection is a competitive process inside the striatum, weighted by dopamine. If you change the weighting --- by depleting dopamine, by downregulating D2 receptors, by chronically saturating D1 sites --- the vote tally changes, and no amount of ``trying'' rewires the tally on the spot.

\textbf{Evolutionary callback.} This is the lamprey's brain. Stephenson-Jones, Samuelsson, Ericsson, Robertson, and Grillner showed in 2011 (\emph{Current Biology}) that all the parts of the mammalian basal ganglia --- striatum, GPi, GPe, STN --- are present in lamprey forebrain, \href{https://pubmed.ncbi.nlm.nih.gov/21700460/?dopt=Abstract}{PubMed} with the same dual D1/D2 direct-and-indirect pathway architecture and the same dopaminergic modulation. Their phrase is worth memorizing: this circuit has been ``conserved, most likely as a mechanism for action selection used by all vertebrates, for over 560 million years.'' \href{https://pubmed.ncbi.nlm.nih.gov/21700460/?dopt=Abstract}{PubMed} Grillner and Robertson (\emph{Current Biology} 2016, ``The Basal Ganglia Over 500 Million Years'') and Grillner, Robertson and Stephenson-Jones (\emph{Journal of Physiology} 2013) drive the point home: the lamprey has the whole thing in miniature, fewer cells, same logic. When I tell you that dopamine is ``ancient,'' I mean that \emph{literally this same wiring diagram} is what a creature with no jaw and no cortex was using to decide whether to swim or stop.

\textbf{Engineering analogy (carefully).} Treat the basal ganglia as a \textbf{gating mechanism} or \textbf{action-selection arbiter}: many candidate actions presented in parallel, one winner unmasked, the rest suppressed. Dopamine is the gain knob on that arbiter. This is the analogy that motivates the gating role of similar structures in models of working memory (O'Reilly and Frank's PBWM, for instance). The analogy is useful as long as you do \emph{not} slip into ``the basal ganglia is a softmax.'' It is messier: it is a leaky, plastic, tonically inhibitory selector with a slow learning rule wired into it.

\textbf{Misconceptions to scrub.} - The basal ganglia is \emph{not} ``the movement part of the brain'' in the way a textbook from 1970 would have it. It selects all kinds of actions --- including cognitive ones, like which thought to pursue. Habit, rumination, procrastination, and ``frozen'' indecision all live here too. - Direct/indirect is not literally ``go and stop.'' Both pathways are active during normal action; they balance each other. What matters is the \emph{ratio}. - ``Dopamine excites everything'' is wrong. D1 neurons are excited; D2 neurons are inhibited. \href{https://pmc.ncbi.nlm.nih.gov/articles/PMC3853485/}{PubMed Central} Same molecule, opposite signs, depending on receptor.

\begin{center}\rule{0.5\linewidth}{0.5pt}\end{center}

\subsection{2. The Striatum --- the input port where everything meets}\label{06_neurology_101-the-striatum-the-input-port-where-everything-meets}

\textbf{Where it sits.} The striatum \emph{is} the input layer of the basal ganglia. Anatomically split into:

\begin{itemize}
\tightlist
\item
  \textbf{Dorsal striatum:} caudate and putamen. Receives projections largely from sensorimotor and associative cortex. Subdivides further into \textbf{dorsomedial striatum (DMS)} --- goal-directed, flexible, ``what do I want and how do I get it'' --- and \textbf{dorsolateral striatum (DLS)} --- habits, automatized sequences, ``I have done this 10,000 times.'' \href{https://www.biorxiv.org/content/10.1101/2021.04.03.438314.full.pdf}{biorxiv}
\item
  \textbf{Ventral striatum:} nucleus accumbens (core and shell) and olfactory tubercle. \href{https://www.pnas.org/doi/10.1073/pnas.1921007117}{PNAS} Receives projections from limbic cortex (orbitofrontal, anterior cingulate, ventromedial PFC), amygdala, and hippocampus. This is the motivational engine --- Pavlovian cues, incentive salience, ``is this thing worth approaching.''
\end{itemize}

\textbf{What it does.} Every cortical area that wants to influence action funnels its vote through here. Cortex glutamate excites medium spiny neurons; dopamine modulates how strongly those cortical votes count and which way they push (D1 vs D2). The striatum is therefore \emph{the} convergence point: cortex (what is happening, what could happen), thalamus (current state), amygdala (does it matter), hippocampus (have I been here before), VTA/SNc (is it worth it).

There is a gradient: ventromedial → dorsolateral striatum maps roughly onto limbic → associative → sensorimotor processing, \href{https://pmc.ncbi.nlm.nih.gov/articles/PMC4240773/}{PubMed Central} and onto motivation → goal-directed action → habit. Habit formation, in animal work, looks like activity migrating along this axis with overtraining (Balleine, O'Doherty, and others; Belin and Everitt's work on the ventral-to-dorsal ``devolution'' in addiction).

\textbf{Connection to the book's argument.} The dorsal/ventral split is going to do a lot of work for us later. The compulsive, ``I don't even want it anymore but I am still doing it'' quality of late-stage stimulant use is partly a story about the locus of control sliding from ventral (motivated) to dorsal (habituated) striatum, \href{https://www.sciencedirect.com/science/article/pii/S0149763413000468}{ScienceDirect} while the dopamine signal that originally trained the habit becomes anhedonic background noise. For now, just keep the geography: ventral = wanting, dorsal = doing. They are continuous, not separate.

\textbf{Evolutionary callback.} Lamprey striatum, with its D1 and D2 MSN populations, has the same input role: cortex/pallium and thalamus glutamate in, dopamine modulating, GABAergic projection to pallidum out. The ventral/dorsal differentiation as we know it expanded enormously in mammals, but the striatal cell types and microcircuit logic are the same shape they were half a billion years ago.

\textbf{Engineering analogy.} Think of the striatum as a \textbf{massive associative input layer} whose synapses are continuously re-weighted by a third signal (dopamine) that says ``that combination of inputs you just received predicted a better-than-expected outcome --- strengthen it.'' It is the place where reinforcement learning in the brain most concretely happens.

\textbf{Misconceptions.} - ``Nucleus accumbens is the pleasure center.'' It is not. It is heavily involved in motivation, incentive salience, and effort, not in the hedonic experience of pleasure per se. \href{https://pubmed.ncbi.nlm.nih.gov/17225164/}{PubMed} Berridge's wanting/liking distinction lives here. - ``Habit is bad, goal-directed is good.'' Habit is computationally cheap and usually correct; the trouble is when a habit outlives the goal that built it.

\begin{center}\rule{0.5\linewidth}{0.5pt}\end{center}

\subsection{3. The Prefrontal Cortex --- the simulation engine that votes, but does not decide}\label{06_neurology_101-the-prefrontal-cortex-the-simulation-engine-that-votes-but-does-not-decide}

\textbf{Where it sits.} The front of the frontal lobe, anterior to motor and premotor cortex. Subdivisions matter:

\begin{itemize}
\tightlist
\item
  \textbf{Dorsolateral PFC (dlPFC):} working memory, rule maintenance, manipulation of held information.
\item
  \textbf{Ventromedial PFC and orbitofrontal cortex:} value, outcome representation, ``is this still worth it.''
\item
  \textbf{Anterior cingulate cortex (dorsal ACC):} conflict, effort allocation, cost monitoring.
\end{itemize}

\textbf{What it does.} The PFC's job, in Miller and Cohen's now-classic 2001 formulation (\emph{Annual Review of Neuroscience}, ``An Integrative Theory of Prefrontal Cortex Function''), is the \textbf{active maintenance of patterns of activity that represent goals and the means to achieve them}, and the use of those patterns as \textbf{bias signals} to other brain structures. That last word is the one to underline. PFC does not make decisions. PFC makes biases. It holds the goal online --- \emph{I am writing this paragraph, not checking my phone} --- and biases the action-selection competition in striatum so that the goal-relevant program gets a slightly bigger vote than the reflex.

This is why the PFC has been called a ``simulation engine'': it can hold representations of states that are \emph{not} currently true (a future, a counterfactual, a rule), and use those representations to push the rest of the brain. It extends the basal ganglia's time horizon from ``what should I do in the next 200 milliseconds based on what is in front of me'' to ``what should I do over the next hour, week, life.''

\textbf{Connection to the book's argument.} ``Willpower'' is a folk-psychology label for what happens when PFC bias is strong enough to overcome a basal-ganglia preference for the easier program. But that bias has a cost; it depends on tonic dopamine in PFC (D1-mediated persistent activity in the dlPFC working-memory circuits --- see Goldman-Rakic, Arnsten); it depends on the goal representation actually being maintained; and it depends on the striatum being in a state where extra biasing input can flip the winner. When dopaminergic tone is low, the bias signal is weaker, the goal representation degrades, and the striatum tilts toward defaults (habit, reflex, the easier option). The subjective experience is ``I just couldn't make myself do it.'' The mechanical reality is: PFC biases were not strong enough to win the vote. There is no separate organ of will.

\textbf{Evolutionary callback.} Lampreys have a pallium with motor, somatosensory, and visual areas (Suryanarayana, Robertson, Wallén, Grillner --- there is a small lamprey ``cortex'' with retinotopic maps). What they do not have is the massively expanded frontal cortex of mammals, particularly the granular dorsolateral PFC of primates. The basal-ganglia loop is ancient; the long simulation engine bolted onto it is recent. So when you feel like there is a tug-of-war between ``the part of me that wants to plan'' and ``the part of me that just wants to act now,'' you are not wrong --- you are feeling 500 million years of evolutionary stratification.

\textbf{Engineering analogy.} Treat the PFC as a \textbf{working-memory buffer holding the current task representation}, plus a \textbf{policy prior} that biases an underlying action-selection module (the basal ganglia). It is roughly the ``context'' in a controller architecture, or the system prompt that conditions everything downstream. It does not output the action; it shapes the distribution from which the action is sampled.

\textbf{Misconceptions.} - ``PFC is the rational part / executive / CEO.'' It is not a homunculus. It is a slow, expensive, dopamine-hungry maintenance system whose output is biasing influence, not commands. - ``PFC stores long-term memories.'' It does not. It holds things online. Long-term storage is hippocampus/cortex consolidation; PFC is the scratchpad. - ``If I just had more PFC, I would have more willpower.'' The bottleneck is usually not PFC capacity but the dopaminergic state that allows PFC to win the bias war.

\begin{center}\rule{0.5\linewidth}{0.5pt}\end{center}

\subsection{4. The Thalamus --- the relay that closes every loop}\label{06_neurology_101-the-thalamus-the-relay-that-closes-every-loop}

\textbf{Where it sits.} Two egg-shaped masses dead-center in the brain, one per hemisphere, sitting on top of the brainstem and below the cortex. Lots of nuclei. The ones that matter for the cortico-basal-ganglia loop are the \textbf{mediodorsal (MD)} and \textbf{ventral anterior / ventral lateral (VA/VL)} nuclei, which receive basal ganglia output and project to PFC and motor cortex respectively. \href{https://www.jneurosci.org/content/27/31/8161}{Journal of Neuroscience}

\textbf{What it does.} Almost every signal that reaches cortex passes through thalamus, and almost every cortical loop closes through thalamus. For our purposes: when the basal ganglia release a thalamic relay from inhibition, that relay activates a cortical area, which then decides what to do next and feeds another vote into striatum. That is how the loop \emph{loops}. Without thalamus, the cortex and basal ganglia are open wires that never close into a circuit.

The thalamus is also where attention modulates traffic --- pulvinar and reticular thalamic nucleus gate which channels cortex actually receives. So ``paying attention'' and ``selecting an action'' are not unrelated; they share thalamic machinery.

\textbf{Connection to the book's argument.} I will not dwell on thalamus much later, because the action is upstream and downstream of it. But its role as the closing relay matters because it explains why basal-ganglia output is sufficient to gate cortex: BG inhibits or releases thalamus, thalamus drives cortex, cortex drives the next round. Disruptions anywhere in this loop --- Parkinson's, Tourette's, OCD, dystonia, drug-induced state changes --- manifest at the thalamic gate.

\textbf{Evolutionary callback.} Lampreys already have a thalamus that receives basal-ganglia output and projects back to pallium (Grillner's group has traced this). The cortico-striato-thalamo-cortical loop is not a primate invention. It is the vertebrate template.

\textbf{Engineering analogy.} A \textbf{switching fabric} or \textbf{bus}, with the basal ganglia's output as one of the gating signals. Useful, but do not push it; thalamus also does substantial computation, not pure relay.

\textbf{Misconceptions.} ``Thalamus is just a relay'' is the line in textbooks; it is approximately true and importantly false. There is genuine processing --- especially in higher-order thalamic nuclei like pulvinar and MD --- but for our chapter, ``loop-closer and gatekeeper'' is enough.

\begin{center}\rule{0.5\linewidth}{0.5pt}\end{center}

\subsection{\texorpdfstring{5. Limbic Structures --- what tells the loop \emph{what matters}}{5. Limbic Structures --- what tells the loop what matters}}\label{06_neurology_101-limbic-structures-what-tells-the-loop-what-matters}

The basal-ganglia/PFC/thalamus loop selects actions. But what biases the selection toward food when you are hungry, away from threat when you are scared, toward the place you remember finding water? Limbic inputs. Three to know.

\subsubsection{Amygdala}\label{06_neurology_101-amygdala}

\textbf{Where.} Almond-shaped cluster in the medial temporal lobe, anterior to hippocampus.

\textbf{What it does.} Detects motivationally significant stimuli --- both threatening and rewarding --- and sends that signal to ventral striatum, hypothalamus, brainstem, and PFC. The basolateral amygdala learns stimulus--outcome associations (this cue predicts that thing). The central amygdala is a major output station that recruits autonomic and behavioral responses and, crucially, projects into mesolimbic circuits.

\textbf{The big correction.} ``Amygdala is the fear center'' is folk neuroscience. Recent work (e.g., the central-amygdala valence/salience papers by Yang and others, 2024--2025; Pignatelli \& Beyeler 2019; Fadok and colleagues) has made it unambiguous that the amygdala --- including the central amygdala once thought to be purely defensive --- encodes both \textbf{valence} (good vs bad) and \textbf{salience} (how much this stimulus matters, regardless of sign). It is recruited by reward as much as by threat. \href{https://www.ncbi.nlm.nih.gov/pmc/articles/PMC11723578/}{nih} Treat the amygdala as a \textbf{valence-and-salience tagger} whose output biases what the striatum considers worth pursuing or avoiding.

\textbf{Book connection.} Anhedonia and the flattening of motivational salience that come with chronic dopamine dysregulation are not just a striatal story; they involve degraded amygdala tagging and degraded amygdala-to-ventral-striatum traffic. When a stimulus stops ``popping,'' part of the popping was happening here.

\subsubsection{Hippocampus}\label{06_neurology_101-hippocampus}

\textbf{Where.} Curled structure in the medial temporal lobe, just behind/below the amygdala.

\textbf{What it does.} Spatial and episodic memory. Where am I, where was I, what happened last time I was here, what is this context. Projects to ventral striatum (especially nucleus accumbens) and to PFC.

\textbf{Book connection.} Context-dependence of motivation and craving lives partly in the hippocampal projection to nucleus accumbens. The reason that walking past a particular doorway can crack open a craving you thought was gone is hippocampal context retrieval feeding a ventral striatum that has been trained to expect dopamine in that context.

\subsubsection{Hypothalamus}\label{06_neurology_101-hypothalamus}

\textbf{Where.} Below the thalamus, above the pituitary, at the base of the brain.

\textbf{What it does.} Bodily-need signals: hunger, thirst, temperature, sleep, sex, stress. The hypothalamus is where homeostatic state becomes a signal that the rest of the brain can read. It influences VTA dopamine neurons and ventral striatum, biasing motivation toward whatever the body currently needs.

\textbf{Book connection.} When sleep deprivation makes everything taste duller and dopamine drugs more compelling, that is hypothalamic state leaning on mesolimbic circuits. The hypothalamus is a constant reminder that ``the brain'' is downstream of a body, and dopamine in particular is sensitive to homeostatic context.

\textbf{Engineering analogy (whole limbic block).} Treat the limbic system as a set of \textbf{biasing inputs} to the cortex--basal-ganglia loop: amygdala feeds in \emph{salience and valence}, hippocampus feeds in \emph{context}, hypothalamus feeds in \emph{bodily-need state}. They do not select actions; they shape the priors over which the loop selects.

\textbf{Misconceptions.} - ``There is a limbic system'' as a clean anatomical entity is itself contested (LeDoux has argued forcefully that ``limbic system'' is a vestigial term). Use it as a convenient bin, not a precise carving. - ``Amygdala = fear, hippocampus = memory'' is the cartoon. Both do more.

\begin{center}\rule{0.5\linewidth}{0.5pt}\end{center}

\subsection{6. The VTA and Substantia Nigra --- where dopamine comes from}\label{06_neurology_101-the-vta-and-substantia-nigra-where-dopamine-comes-from}

\textbf{Where they sit.} Both in the midbrain (mesencephalon). - \textbf{Ventral tegmental area (VTA):} medial midbrain. Dopamine neurons project to \textbf{nucleus accumbens, ventral striatum, amygdala, hippocampus, and PFC} (the \textbf{mesolimbic} and \textbf{mesocortical} pathways). - \textbf{Substantia nigra pars compacta (SNc):} more lateral, in the midbrain. Dopamine neurons project to \textbf{dorsal striatum} (the \textbf{nigrostriatal} pathway). Its degeneration is the proximate cause of Parkinson's disease.

(Substantia nigra also has a \emph{pars reticulata}, SNr, which is \emph{not} dopaminergic --- it is a basal-ganglia output station, sister to GPi. Same anatomical neighborhood, completely different job. Easy to confuse.)

\textbf{What they do.} They make dopamine and spray it across striatum and PFC according to specific firing patterns. The patterns matter:

\begin{itemize}
\tightlist
\item
  \textbf{Tonic firing:} slow, baseline. Sets the ambient dopamine level. Influences how much effort feels available, how readily striatum responds to cortical input.
\item
  \textbf{Phasic bursts:} brief, high-frequency. Encode the moment-to-moment salience and prediction-error signals.
\end{itemize}

What does that phasic signal compute? This is one of the questions the book will keep coming back to, because there are at least three serious answers, and they are all partly right:

\begin{itemize}
\item
  \textbf{Reward prediction error (Schultz).} Wolfram Schultz's monkey-recording work, formalized with Dayan and Montague (\emph{Science} 1997), showed that midbrain dopamine neurons fire when reward is \emph{better than expected}, stay quiet when reward is \emph{as expected}, and dip when reward is \emph{worse than expected}. \href{https://www.tandfonline.com/doi/full/10.31887/DCNS.2016.18.1/wschultz}{Taylor \& Francis} This is the temporal-difference error of reinforcement learning, in neurons. (Schultz, ``Dopamine reward prediction error coding,'' \emph{Dialogues in Clinical Neuroscience} 2016, gives the clean review.) This is the dominant computational story for phasic dopamine.
\item
  \textbf{Incentive salience / precision weighting (Berridge; Friston).} Berridge's work distinguishes ``wanting'' (incentive salience, dopamine-dependent) from ``liking'' (hedonic impact, more opioid- and endocannabinoid-dependent). Dopamine makes things \emph{pop}, which is not the same as making them \emph{feel good}. Friston and colleagues (\emph{PLOS Computational Biology} 2012, ``Dopamine, Affordance and Active Inference''; Schwartenbeck et al.~\emph{Cerebral Cortex} 2015; and a 2015 \emph{Frontiers in Computational Neuroscience} paper on dopamine, reward learning, and active inference) recast this in Bayesian terms: dopamine encodes the \textbf{precision} (inverse uncertainty) of beliefs about which actions or cues are worth taking seriously. \href{https://pmc.ncbi.nlm.nih.gov/articles/PMC3252266/}{PubMed Central} Higher dopamine = higher confidence in the current policy, more decisive action, sharper attention to selected cues; lower dopamine = everything is hazy, no policy is precise enough to commit to.
\item
  \textbf{Effort cost / behavioral activation (Salamone).} John Salamone's program of work (Salamone, Correa, Farrar, Nunes, and Pardo, \emph{Frontiers in Behavioral Neuroscience} 2009, ``Dopamine, Behavioral Economics, and Effort''; Salamone \& Correa's reviews; the 2016 \emph{Brain} review on activational and effort-related aspects of motivation) has shown, very cleanly, that depleting nucleus-accumbens dopamine does not abolish liking food and does not abolish the ability to learn --- it specifically reduces willingness to \emph{work} for reward. Animals with low accumbens dopamine still eat freely available food but will not press a lever many times for tastier food; they reallocate to lower-effort alternatives. \href{https://academic.oup.com/brain/article/139/5/1325/2468761}{Oxford Academic} Dopamine, in this account, is the signal that says \emph{the effort is worth it}.
\end{itemize}

These are not three competing theories. They are three computational facets of the same neuromodulator, depending on which aspect of behavior you measure. Phasic dopamine looks like prediction error in Schultz's tasks; tonic accumbens dopamine looks like effort-willingness in Salamone's tasks; and the precision/salience formulation in Friston's framework arguably subsumes both at a higher level of description. For the book's purposes you should hold all three in mind:

\begin{quote}
\textbf{Dopamine is a learning signal, an effort signal, and a precision signal. None of these is ``pleasure.''}
\end{quote}

\textbf{Connection to the book's argument.} This is the single most important paragraph in the chapter. When stimulants chronically saturate dopamine receptors and the brain compensates by downregulating D2/D3 sites, every one of these computational roles degrades simultaneously. Prediction errors are blunted (nothing surprises, nothing teaches). Salience precision collapses (nothing pops). Effort signals fail (nothing feels worth doing). The subjective experience is a fused symptom --- anhedonia, cognitive bradykinesia, ``I cannot make myself care'' --- but mechanistically it is the same signal failing in three different roles at once. Calling that failure a deficit of ``willpower'' is like calling a power outage a deficit of motivation.

\textbf{Evolutionary callback.} Lampreys have midbrain dopamine neurons projecting to striatum, with the same D1/D2 receptor scheme on striatal output cells. Thompson, Ménard, Pombal, and Grillner (2008) showed that dopamine depletion in lamprey impairs motor behavior, just as it does in mammalian Parkinson's. The mesolimbic dopamine system is not a mammalian luxury. It is the original throttle.

You can also push the callback further back, to the chapter we already wrote: the molecular logic of dopamine --- a neuromodulator that gates which inputs count, biasing plasticity at glutamatergic synapses with calcium signaling as the downstream lever --- is the ancient eukaryotic logic of ``this signal matters, attend to it'' routed through a midbrain nucleus. Dopamine is what happens when 3-billion-year-old ``this stimulus is salient'' chemistry gets a dedicated postal service.

\textbf{Engineering analogy.} Tempting to say ``dopamine = reward signal in RL,'' and that is fine for the Schultz half. The fuller analogy is: dopamine is a \textbf{multi-purpose modulator that simultaneously updates value estimates (RPE), sets the gain on action selection (gating), and weights confidence in the current policy (precision)}. In artificial systems we tend to factor these into separate components. The brain does not.

\textbf{Misconceptions to scrub aggressively.} - ``Dopamine is the pleasure chemical.'' No.~It is a learning, salience, and effort modulator. The hedonic gloss came from early misinterpretations of self-stimulation experiments and has been wrong for thirty years. - ``More dopamine = more happy.'' No.~Tonically elevated dopamine in the wrong contexts produces craving, intrusive salience, paranoia, and psychosis, not happiness. - ``VTA = reward, SNc = movement.'' Shorthand at best. SNc dopamine is involved in reward-related learning in dorsal striatum (the nigrostriatal pathway underwrites habit learning and stimulus-response associations); VTA dopamine influences cognition and movement via PFC and accumbens. Same neurotransmitter, different terminals. - ``Drugs of abuse just dump dopamine.'' They do, but the \emph{important} part is what chronic dumping does to receptor density, baseline tone, and the precision of the prediction-error and effort signals. The acute spike is not the pathology; the chronic adaptation is.

\begin{center}\rule{0.5\linewidth}{0.5pt}\end{center}

\subsection{Putting it back together --- the napkin sketch}\label{06_neurology_101-putting-it-back-together-the-napkin-sketch}

Here is what I want you to be able to draw:

\begin{enumerate}
\def\labelenumi{\arabic{enumi}.}
\tightlist
\item
  A box labeled \textbf{PFC} at the top.
\item
  An arrow from PFC down to a box labeled \textbf{Striatum}.
\item
  Two arrows from \textbf{Limbic} (amygdala, hippocampus, hypothalamus) also pointing at Striatum.
\item
  From Striatum, two arrows: a \textbf{direct} (``go,'' D1) and an \textbf{indirect} (``no-go,'' D2) pathway, both ending at a box labeled \textbf{GPi/SNr} (the basal ganglia output).
\item
  From GPi/SNr, an inhibitory arrow to \textbf{Thalamus}.
\item
  From Thalamus, an arrow back up to \textbf{PFC} (and to motor cortex).
\item
  A separate box labeled \textbf{VTA / SNc} with a thick modulatory arrow that touches Striatum, PFC, and the limbic structures.
\end{enumerate}

That is the loop. Read it as a sentence: \emph{``Cortex and limbic inputs send votes to the striatum; the striatum sends a winner through direct/indirect competition to the basal ganglia output; the output gates the thalamus; the thalamus closes the loop back to cortex; dopamine from the midbrain modulates how the votes count and what gets learned.''}

Everything later in this book --- D2/D3 downregulation, the failure of stimulant-replenished motivation to feel like motivation, cognitive bradykinesia, the way ``willpower'' decomposes into effort signals and precision weights and direct/indirect ratios --- is a perturbation of that diagram.

\begin{center}\rule{0.5\linewidth}{0.5pt}\end{center}

\subsection{Why ``willpower'' is a bad term for this system}\label{06_neurology_101-why-willpower-is-a-bad-term-for-this-system}

Let me close on this, because it is the through-line.

``Willpower'' implies a unitary mental quantity that you have more or less of, and that you spend like a currency. Nothing in the picture above corresponds to such a quantity. What actually exists is:

\begin{itemize}
\tightlist
\item
  a \textbf{policy bias} maintained in PFC (which depends on tonic dopamine in PFC, on glucose, on sleep, on whether the goal representation is even being held online);
\item
  a \textbf{competition} in striatum between D1-mediated ``go'' and D2-mediated ``no-go'' votes for many simultaneously available programs;
\item
  a \textbf{modulatory tone} of dopamine from VTA/SNc that determines the gain of that competition and the precision of the bias;
\item
  \textbf{limbic inputs} that constantly push the competition toward whatever currently looks salient;
\item
  a \textbf{thalamic gate} that lets the winner reach motor and cognitive output.
\end{itemize}

When somebody ``fails to exert willpower,'' at least one of those pieces failed. The PFC bias may have been weak because the goal representation decayed. The striatal competition may have been tilted because chronic dopaminergic perturbation downregulated D2 receptors and shifted the direct/indirect balance. The dopaminergic precision signal may have been too low to make the goal-relevant cue feel real. The amygdala may have been screaming about an unrelated threat. None of these is a moral failure. None of them is fixable by ``trying harder,'' because trying harder is itself a function of the same machinery that just failed.

This is not an argument for fatalism. It is an argument for \textbf{specificity}. Once you know which piece failed, you can sometimes fix it --- or at least stop blaming a person for the failure of a circuit they did not design and cannot see. That is the entire pedagogical point of this chapter, and it is what the rest of the book will lean on.

\begin{center}\rule{0.5\linewidth}{0.5pt}\end{center}

\subsection{A short note on sources and uncertainty}\label{06_neurology_101-a-short-note-on-sources-and-uncertainty}

Where I have been confident, I have been confident. Where the field is still arguing, I have flagged it (the three readings of dopamine; the carving of ``limbic system''; the relative roles of DMS/DLS and ventral/dorsal striatum in addiction). The major sources behind this chapter are:

\begin{itemize}
\tightlist
\item
  \textbf{Lamprey basal ganglia conservation:} Stephenson-Jones et al., \emph{Current Biology} 2011; Stephenson-Jones et al., \emph{Journal of Comparative Neurology} 2012; Grillner, Robertson \& Stephenson-Jones, \emph{Journal of Physiology} 2013; Grillner \& Robertson, \emph{Current Biology} 2016 (``The Basal Ganglia Over 500 Million Years''); Robertson et al., \emph{Progress in Brain Research} 2014; Suryanarayana et al.~on lamprey pallium; Thompson, Ménard, Pombal \& Grillner 2008 on dopamine depletion in lamprey.
\item
  \textbf{Dopamine reward prediction error:} Schultz, Dayan \& Montague, \emph{Science} 1997; Schultz, \emph{Dialogues in Clinical Neuroscience} 2016; Schultz, \emph{Nature Reviews Neuroscience} 2016 (two-component response); Schultz, \emph{F1000Research} 2019.
\item
  \textbf{Dopamine as precision / active inference:} Friston, Shiner, FitzGerald et al., \emph{PLOS Computational Biology} 2012 (``Dopamine, Affordance and Active Inference''); Schwartenbeck, FitzGerald, Mathys, Dolan \& Friston, \emph{Cerebral Cortex} 2015; FitzGerald, Dolan \& Friston, \emph{Frontiers in Computational Neuroscience} 2015.
\item
  \textbf{Dopamine and effort:} Salamone, Correa, Farrar, Nunes \& Pardo, \emph{Frontiers in Behavioral Neuroscience} 2009 (``Dopamine, Behavioral Economics, and Effort''); Salamone \& Correa, \emph{Neuron} 2012; Salamone et al., \emph{Brain} 2016 (activational/effort-related aspects of motivation); Salamone \& Correa, \emph{Annual Review of Psychology} 2024/2025 (neurobiology of activational motivation).
\item
  \textbf{PFC function:} Miller \& Cohen, \emph{Annual Review of Neuroscience} 2001 (``An Integrative Theory of Prefrontal Cortex Function''); Goldman-Rakic and Arnsten on dopamine and dlPFC working memory.
\item
  \textbf{Striatum dorsal/ventral:} Balleine \& O'Doherty, \emph{Neuropsychopharmacology} 2010; Belin \& Everitt on the ventral-to-dorsal ``devolution'' of drug seeking; Voorn, Vanderschuren, Groenewegen, Robbins \& Pennartz, \emph{Trends in Neurosciences} 2004 (``Putting a spin on the dorsal-ventral divide of the striatum'').
\item
  \textbf{Amygdala valence/salience:} Pignatelli \& Beyeler, \emph{Current Opinion in Behavioral Sciences} 2019; Fadok et al.~work on central amygdala circuits; the recent eLife paper on valence and salience encoding in central amygdala (2024--2025).
\item
  \textbf{CSTC loops generally:} Alexander, DeLong \& Strick's 1986 \emph{Annual Review of Neuroscience} paper on parallel cortico-basal-ganglia-thalamo-cortical loops remains the foundational reference.
\end{itemize}

I have not made up any citations. Where I summarize I am summarizing; where I quote, the source is named in line. If a claim above feels too clean, treat it as a teaching simplification of a messier literature, and assume that the messier literature exists.

That is enough neuroanatomy. Next chapter we point this loop at the actual problem.
